\section{Model Distillation} \label{sec:23-background-modeldistill}
\subsection{The Original Formulation} \label{sec:23-background-modeldistill-hinton}
\begin{foldable}
  The idea of training a smaller student network to mimic a larger teacher originated with model compression~\cite{bucila2006model}.
  This was formalized as \textit{knowledge distillation} (\ac{KD})~\cite{hinton2015distilling}, proposing that students learn not from hard labels but from soft probabilistic targets that encode the teacher's learned structure.
  Model distillation trains a compact \textit{student} network $S$ to replicate the behaviour of a large, pre-trained \textit{teacher} network $T$, illustrated in Figure~\ref{fig:23-background-modeldistill-overview}.
  The goal is to transfer the teacher's learned \textit{decision-boundary structure} and relational knowledge, the structured understanding of the task and its geometry acquired over the full training budget, rather than requiring pointwise input-output fidelity on a fixed set of training examples.
  This distinguishes distillation from model approximation, which targets pointwise structural matching, and from label smoothing, which perturbs labels uniformly without reference to the teacher's learned geometry.
\end{foldable}

\begin{figure}[ht]
  \centering
  \includegraphics[width=0.80\linewidth]{./figures/20-background/23-background-modeldistill-overview.pdf}
  \caption[The original knowledge distillation formulation.]{
    The original knowledge distillation formulation: a student network is trained on the same input data as the teacher, learning to match the teacher's soft output distribution in addition to the hard labels alone.
  }
  \label{fig:23-background-modeldistill-overview}
\end{figure}

\begin{foldable}
  The obstacle is that standard supervised training discards the very information distillation needs.
  A one-hot label records which class won the annotation decision and assigns zero probability to every other class.
  A large network trained on such labels does not merely learn to classify correctly; it implicitly acquires a structured posterior over the class space, shaped by the statistical regularities of the full training distribution.
  When such a network evaluates an image of a cat, it assigns high probability to \texttt{cat}, a smaller amount to \texttt{dog}, and a trace to \texttt{leopard}, while assigning essentially nothing to \texttt{car}, encoding which classes share visual structure and which do not.
  Similarly, a digit \texttt{8} receives non-trivial mass on \texttt{3} and near-zero on \texttt{1}, reflecting the perceptual similarity the training data induced.
  This inter-class relational structure, termed \textit{dark knowledge}~\cite{hinton2015distilling}, is present in every forward pass of the teacher yet discarded by hard-label supervision.
\end{foldable}

\begin{foldable}
  \Acl{KD} recovers dark knowledge by training the student to minimise the \ac{KL} divergence between the teacher's output distribution $q$ and the student's $p$:
  \begin{equation}
    \mathcal{L}_{\text{KD}} = \mathcal{L}_\text{KL}(q \| p).
    \label{eq:kd-loss-core}
  \end{equation}
  In practice, a highly confident teacher concentrates nearly all mass on the correct class, making dark knowledge numerically invisible even though it exists.
  \textit{Temperature scaling} addresses this: a parameter $\tau \geq 1$ applied to the pre-softmax logits softens the distribution and amplifies inter-class structure:
  \begin{equation}
    q_k^{(\tau)} = \frac{\exp(z_k^{T} / \tau)}{\sum_{j=1}^{K} \exp(z_j^{T} / \tau)}.
    \label{eq:softmax-temperature}
  \end{equation}
  A hard-label cross-entropy term is added to keep the student anchored to ground truth, yielding the full objective~\cite{hinton2015distilling}:
  \begin{equation}
    \mathcal{L}_{\text{KD}}
    = \alpha \tau^{2} \mathcal{L}_\text{KL}(q^{(\tau)} \| p^{(\tau)})
    + (1 - \alpha) \mathcal{L}_{\text{CE}}(p^{(1)}, y),
    \label{eq:kd-loss}
  \end{equation}
  where $\alpha \in [0,1]$ balances the two terms and $\tau^{2}$ restores gradient scale after softening.
\end{foldable}

\begin{foldable}
  The teacher's output is not merely a prediction: it is a compressed summary of everything the teacher learned from the full training distribution, encoded in a probability vector.
  Training the student against this posterior transfers the teacher's generalization directly: its responses on hard examples, its calibrated uncertainty, and its implicit encoding of class relationships.
  Empirically, a student of identical architecture to the teacher consistently outperforms the same model trained on hard labels alone~\cite{hinton2015distilling}, confirming that the gain comes from signal quality rather than capacity.
  Subsequent work has shown that the teacher's intermediate representations and the relational structure among samples carry additional signal, motivating a broader taxonomy of distillation methods.
\end{foldable}

\subsection{Taxonomy of Model Distillation Methods} \label{sec:23-background-modeldistill-taxonomy}
\begin{foldable}
  The original formulation transfers knowledge through the teacher's output distribution, but the field has since expanded to exploit richer sources of teacher signal.
  The design space of \ac{KD} decomposes along two orthogonal dimensions: (1) the \textit{source of knowledge} transferred (response-based, feature-based, or relation-based methods), and (2) the \textit{training scheme} that structures the teacher-student relationship (offline, online, self, or co-distillation).
  We first inventory the knowledge sources by examining what each family transfers and what teacher access is required, then discuss training schemes.
  Together they characterise the full design space.
\end{foldable}

\begin{figure}[ht]
  \centering
  \includegraphics[width=0.70\linewidth]{./figures/20-background/23-background-modeldistill-responsekd.pdf}
  \caption[Response-based knowledge distillation.]{
    Response-based distillation matches the teacher's output distribution, requiring only output-level access.
  }
  \label{fig:23-background-modeldistill-responsekd}
\end{figure}

\begin{foldable}
  \textbf{Response-based distillation}~\cite{hinton2015distilling} trains the student to match the teacher's output class distribution:
  \begin{equation}
    \mathcal{L}_{\text{KD}}^{\text{response}} = \alpha \tau^2 \mathcal{L}_{\text{KL}}(q^{(\tau)} \| p^{(\tau)}) + (1-\alpha) \mathcal{L}_{\text{CE}}(p, y),
    \label{eq:kd-response}
  \end{equation}
  as formalized in Equation~\ref{eq:kd-loss}.
  Variants include regression on the raw un-softened logit vector and born-again networks, in which a same-architecture student is sequentially re-distilled from the previous generation, yielding ensembles that outperform the original teacher~\cite{ba2014deep}.
  Response-based methods require only the teacher's output probability vector and are therefore compatible with any setting in which that vector is accessible, including black-box \acp{API}.
\end{foldable}

\begin{figure}[ht]
  \centering
  \includegraphics[width=0.65\linewidth]{./figures/20-background/23-background-modeldistill-featurekd.pdf}
  \caption[Feature-based knowledge distillation.]{
    Feature-based distillation matches intermediate layer representations via a learned projector, requiring internal teacher access.
  }
  \label{fig:23-background-modeldistill-featurekd}
\end{figure}

\begin{foldable}
  \textbf{Feature-based distillation} transfers knowledge from the teacher's intermediate representations by matching aligned feature maps at selected layers:
  \begin{equation}
    \mathcal{L}_{\text{KD}}^{\text{feature}} = \left\| W f^T - f^S \right\|^2,
    \label{eq:kd-feature}
  \end{equation}
  where $f^T$ and $f^S$ are teacher and student feature maps at a chosen intermediate layer, and $W$ is a learned projection matrix.
  This loss can be summed over multiple selected layer pairs.
  FitNets~\cite{romero2015fitnets} introduced matching intermediate hidden layers via a thin regressor, enabling students deeper and thinner than the teacher.
  Attention Transfer~\cite{zagoruyko2016paying} distils spatial attention maps, capturing \textit{where} the teacher focuses rather than \textit{what} it predicts.
  PKT~\cite{passalis2018learning} matches kernel density estimates of teacher and student feature distributions, encoding distributional geometry rather than individual activation values.
  Feature-based methods consistently improve on response-based baselines, but their incompatibility with opaque \acp{API} is fundamental.
  Without access to the teacher's activation dimensionality and representation geometry, the projector bridging teacher and student spaces is entirely ill-posed: there is no anchor in the teacher's representation space to which the student can be steered.
\end{foldable}

\begin{figure}[ht]
  \centering
  \includegraphics[width=0.85\linewidth]{./figures/20-background/23-background-modeldistill-relationkd.pdf}
  \caption[Relation-based knowledge distillation.]{
    Relation-based distillation preserves pairwise distance and angle relationships between samples in the teacher's embedding space.
  }
  \label{fig:23-background-modeldistill-relationkd}
\end{figure}

\begin{foldable}
  \textbf{Relation-based distillation} transfers higher-order structure: pairwise or group-wise relationships among samples in the teacher's embedding space:
  \begin{equation}
    \mathcal{L}_{\text{KD}}^{\text{rel}} = \sum_{i,j} \left[ (d_T(f_i^T, f_j^T) - d_S(f_i^S, f_j^S))^2 + (\Delta\theta_T - \Delta\theta_S)^2 \right],
    \label{eq:kd-relation}
  \end{equation}
  where $f_i^T$ and $f_i^S$ are feature vectors for samples $i$ and $j$ extracted from an intermediate layer, $d$ measures pairwise Euclidean distance, and $\Delta\theta$ is the angle between feature vectors, with subscripts $T$ and $S$ denoting teacher and student respectively.
  RKD~\cite{park2019relational} penalizes discrepancies in distance and angle relationships between sample pairs, arguing that relational information is more transferable across architectures than pointwise feature values.
  Contrastive distillation methods extend this by using self-supervised objectives to preserve the teacher's instance discrimination structure.
  Relation-based methods therefore share the access problem of feature-based methods, as pairwise statistics require intermediate-level teacher outputs.
  Beyond access, a student trained solely on hard labels develops an embedding geometry incommensurable with the teacher's, making the relational target unreachable.
\end{foldable}

\begin{foldable}
  The second axis organises methods by the \textit{training scheme}: the relationship between teacher and student during training, illustrated in Figure~\ref{fig:23-background-modeldistill-schemes}.
  \textit{Offline distillation} uses a fixed, pre-trained teacher queried through its outputs; the student is trained entirely from the teacher's frozen predictions.
  \textit{Online distillation} trains teacher and student simultaneously, allowing the teacher to evolve during the student's training.
  \textit{Self-distillation} requires no separate teacher: the student distils knowledge from its own earlier or shallower representations.
  \textit{Co-distillation} extends this to an ensemble of students that mutually supervise one another.
  All work in this thesis operates in the offline regime.
\end{foldable}

\begin{figure}[ht]
  \centering
  \includegraphics[width=0.99\linewidth]{./figures/20-background/23-background-modeldistill-schemes.pdf}
  \caption[Knowledge distillation training schemes.]{
    Knowledge distillation training schemes: offline (frozen pre-trained teacher), online (jointly trained), self-distillation, and co-distillation.
  }
  \label{fig:23-background-modeldistill-schemes}
\end{figure}

\begin{foldable}
  Across all three source-of-knowledge families, two conditions are implicit: the full training dataset $\mathcal{D}$ is available to supply inputs to the teacher, and the teacher exposes at least its output probability vector.
  When either condition fails, the picture changes dramatically.
  Response-based methods' black-box compatibility becomes their sole advantage when these assumptions break: feature-based and relation-based methods become inaccessible, while response-based distillation survives, albeit degraded, under hard-label-only settings.
  The next subsection formalizes these two assumptions and characterises what happens when they fail.
\end{foldable}

\subsection{Standard Assumptions and Their Limitations} \label{sec:23-background-modeldistill-assumptions}
\begin{foldable}
  \textbf{Assumption 1: the full training dataset is available.}
  In many realistic settings, the original training data is inaccessible: it may contain private records (medical imaging, financial transactions), it may be proprietary to the model provider, or it may have been discarded after training to reduce storage costs.
  Data regulations such as \ac{GDPR}~\cite{eu2016gdpr} or \ac{HIPAA}~\cite{us1996hipaa} further restrict the retention and transfer of training sets across organizational boundaries.
  When $\mathcal{D}$ is unavailable, the student must be trained on a surrogate: a small held-out subset, synthetic data from a generative model, or public proxy data from a related domain.
  This constraint exists on a spectrum, from \textit{few-shot} settings in which only a handful of samples are available to fully \textit{data-free} settings in which no original data exists; the data-free end is strictly harder, as no real distribution signal remains.
\end{foldable}

\begin{foldable}
  \textbf{Assumption 2: the teacher is accessible at the logit level or deeper.}
  All three \ac{KD} families require at minimum the teacher's output probability vector over all $K$ classes; feature-based and relation-based methods additionally require internal activations.
  Modern production models are routinely deployed behind \acp{API} that return only a top-1 predicted label, or with associated confidence scores, to protect proprietary model internals and limit adversarial probing.
  In this setting, the access restrictions eliminate two of the three families outright: feature-based and relation-based methods are impossible without internal activations.
  Response-based methods remain the sole viable family, though they are severely degraded by the loss of soft probability distributions: a single hard label carries $\log_2 K$ bits per query, whereas the teacher's full soft target encodes the inter-class similarity structure it learned.
\end{foldable}

\begin{foldable}
  Together, these two failures characterise the landscape of practical \ac{KD}.
  Absent original data, the distillation signal loses coverage of the true input distribution; absent soft targets, it loses the relational structure that distinguishes \ac{KD} from ordinary supervised training.
  \S\ref{sec:24-background-kdrestricted} surveys these two restriction axes, frames distributional fidelity as the binding constraint, and identifies the structural gaps that motivate Chapters~\ref{ch:30-research01}--\ref{ch:50-research03}.
\end{foldable}
